\begin{abstract}
Transformer layers evolve representations through learned additive
transformations that either preserve the current direction or introduce
direction-changing content. We decompose each transformation into a parallel
component and a perpendicular component, applying the same geometry in
residual space and in attention value space. Across pretrained models, learned
attention and MLP outputs contain substantial parallel components. These
components are not uniformly redundant: inference-time scaling shows that
robustness depends on which space is manipulated. Manipulating the part of
attention aggregation aligned with the projected self-value remains close to
baseline, whereas residual-space scaling is more fragile and perpendicular
scaling is more disruptive. We explain this asymmetry with an attention-side
factorization that separates diagonal self-routing, value self-alignment, and
cross-token aggregation, showing why the projected self-value is the effective
direction for stable value-space edits. The same geometry also supports two
downstream uses. For compression, stronger compressed models have lower
perpendicular error in layer transformations, while parallel error is less
discriminative. For training, retained pretraining runs indicate that
suppressing self-value-parallel components shifts optimization toward more
direction-changing updates and improves average downstream scores. Taken
together, this geometry is not only descriptive, but also a practical
principle for transformer editing, compression diagnostics, and training-time
intervention.
\end{abstract}
